\uuid{M3po}
\exo7id{4072}
\titre{exo7 4072}
\auteur{quercia}
\organisation{exo7}
\datecreate{2010-03-11}
\isIndication{false}
\isCorrection{true}
\chapitre{Equation différentielle}
\sousChapitre{Equations différentielles linéaires}
\cpge{mpsi-mp2i.applications-lineaires.a, pcsi.applications-lineaires.a, ptsi.applications-lineaires.a, mpsi-mp2i.applications-lineaires.b, pcsi.applications-lineaires.b, ptsi.applications-lineaires.b, bcpst1.applications-lineaires-matrices, mp-mpi.reduction-endomorphismes.b, pc.reduction.a, psi.reduction-endomorphismes.a, pt.reduction-endomorphismes.a, bcpst2.valeurs-propres-vecteurs-propres.a}
\module{Analyse}
\niveau{L2}
\difficulte{}

\contenu{
\texte{
Soit $A$ un polynôme à coefficients réels de degré 2 donné.
Au polynôme $P$ de degré inférieur ou égal à $2n$ on fait correspondre le
polynôme $Q = AP' - nA'P$.
}
\begin{enumerate}
    \item \question{Montrer qu'on définit ainsi un endomorphisme $\Phi$ de $\R_{2n}[X]$.}
    \item \question{Chercher les valeurs propres et les vecteurs propres de $\Phi$ dans les
    cas particuliers :\par
   }
    \item \question{$A = X^2-1$,}
    \item \question{$A = X^2$,}
    \item \question{$A = X^2+1$.}
\reponse{
$\lambda = 2k$, $P = \alpha(X-1)^{n-k}(X+1)^{n+k}$ pour $-n \le k \le n$.
$\lambda = 0$, $P = \alpha X^{2n}$.
$\lambda = 0$, $P = \alpha (X^2+1)^n$.
}
\end{enumerate}
}
